% UNCONFIRMED SUPPLEMENT DRAFT (agent, 2026-11-01). Goes with
% main_REFRAME_UNCONFIRMED.tex. Text marked "moved verbatim" is the old main
% text (main.tex at 3b1a24c) unchanged. S1 rows follow the test order of
% EXTENSIONS_TEST_REPORT.md. The printed text names no test numbers.
\documentclass[10pt]{article}
\usepackage[margin=1in]{geometry}
\usepackage{booktabs}
\usepackage{amsmath}
\usepackage{xcolor}
\newcommand{\pending}[1]{\textcolor{red}{[PENDING: #1]}}

\title{Supplement: Equal Pixels, Unequal Views}
\date{}
\begin{document}
\maketitle

\section*{S1. Every registered test and where it is reported}

All tests were registered in the pre-registration or one of its amendments
before any of their numbers existed. ``Main'' gives the section of the paper.
Fit split means \texttt{ntire\_val\_hard}, on which squish was chosen.

\begin{center}
\small
\begin{tabular}{p{6.0cm}p{4.6cm}p{3.4cm}}
\toprule
question & outcome & reported \\
\midrule
proxy: squish-trained beats crop-trained, 3 seeds & passed, +0.078 to +0.093 & Main, Training \\
full fused models on the hard split & passed, +0.074 & Main, Training \\
each external benchmark & SocialRF pass, CommunityAI confounded, Synthbuster fail & Main, Held-out \\
CommunityAI counts only if size-matched JPEG gain clears zero & did not clear, confounded & Main, Held-out \\
aspect-preserving presentations keep 75\% of the gain & passed, 93 to 100\% & Main (one line), S2 \\
squish helps public detectors & UnivFD at chance, B-Free fails & Main, Held-out \\
CLIP repeat of the proxy comparison and test-time gain & passed & Main, Training \\
does squish's frequency content explain the gain? & coverage holds, design confounded & S8 \\
does the gain grow with degradation level? & passed, rise holds & Main, What the gain is made of; S9 \\
field of view or shrink ratio, by regression & not identifiable, no result & S8 \\
is the gain spread across degradation families? & JPEG gains more, mostly headroom & S8 \\
is training field of view narrow? & no, 0.391 to 1.000 after augmentation & Main, Table 1 \\
routed fusion keeps the gain and removes the clean cost & failed & Main, clean cost (one line) \\
multi-scale training matches squish & squish ahead, gap 73--81\% closed & Main, Training \\
crop baseline handicapped at 224? & no & Main, Training \\
learned combiner over nine windows matches squish & squish ahead, B ties on small images & Main, Training \\
does the periphery carry the evidence? & failed, wrong geometry & S6 \\
where is the clean cost on large images? & not located & Main (one line), S7 \\
wording corrections with numbers & no condition & applied in text \\
does the degradation rise survive degradation after squish? & A 86\% survives, B about half & S3 \\
coverage against downscale, NTIRE & both effects & Main, What the gain is made of \\
native 512 window for branch A & +0.042, mixed with view & Main, Training \\
coverage against downscale, SocialRF & downscale only & Main, What the gain is made of \\
downscale by 2: degradation removal or field of view? & field of view & Main, What the gain is made of \\
degradation removal at stronger downscales & larger at 4x than 2x on both branches, not proof & Main, What the gain is made of \\
enlargement at equal field of view & +0.011, interval includes zero, unresolved & Main, Training \\
coverage against downscale on CLIP & field of view, small removal & Main, What the gain is made of \\
downscale on B-Free & +7.9 points on large degraded SocialRF & Main, Held-out \\
\bottomrule
\end{tabular}
\end{center}

Not run: a third test read, once planned for routed fusion, was not run by decision.
Exploratory, not registered: the DINOv2 test-time comparison on the proxies
(S4), the CLIP trained-minus-test-time comparison, and a headroom check on
the degradation-family test (S8).

\section*{S2. Resolution-shortcut details (moved verbatim)}

\paragraph{The gain is not only the geometric distortions.} The challenge's
validation splits add random crops, aspect crops and downscaling, and the
hard validation and test splits also add perspective transforms; its training
split has none of these.
These change native size and could carry the effect. Grouping the distorted
rows of \texttt{ntire\_val\_hard} by their transforms
(Table~\ref{tab:groups}), squish gains in every group. On rows that received
only photometric transforms the gain is +9.0 points on A and +8.3 on B, with
intervals clear of zero. Those rows have a median short side of 784 pixels;
large native images are enough to produce the effect.

\begin{table}[h]
\centering
\small
\begin{tabular}{lrrcc}
\toprule
distortions & n & short p50 & A gain & B gain \\
\midrule
photometric only & 241 & 784 & +0.090 & +0.083 \\
downscale, no crop & 740 & 432 & +0.095 & +0.065 \\
crop or aspect crop & 191 & 320 & +0.059 & +0.063 \\
perspective & 78 & 536 & +0.146 & +0.140 \\
\bottomrule
\end{tabular}
\caption{Squish gain in robust AUC over the centre window, within groups of
distorted \texttt{ntire\_val\_hard} rows, frozen crop-trained branches. All
eight 95\% intervals exclude zero.}
\label{tab:groups}
\end{table}

\paragraph{Aspect ratio.} Squish stretches non-square images, which
Grommelt et al.\ warn introduces frequency artefacts
along one axis. To attribute the gain to field of view
rather than to the stretch, we registered two aspect-preserving presentations
before scoring them: \emph{short-centre} (short side to the input size, then
the centre square) and \emph{letterbox} (long side to the input size, padded
to a square). We predicted that each would recover at least 75\% of squish's
gain on both branches. Both did (Table~\ref{tab:aspect}): short-centre kept
93\% of the gain on A and 97\% on B, letterbox 100\% and 94\%. Against squish,
short-centre is 0.7 points lower in robust AUC on A, with an interval that
excludes zero, and indistinguishable on B. Letterbox matches squish in robust
AUC but loses 0.4 points of clean AUC on both branches. The stretch is
therefore not what drives the gain; showing the whole short side is.

\begin{table}[h]
\centering
\small
\begin{tabular}{lcccc}
\toprule
 & \multicolumn{2}{c}{robust gain over centre} & \multicolumn{2}{c}{clean vs.\ squish} \\
\cmidrule(lr){2-3}\cmidrule(lr){4-5}
presentation & A & B & A & B \\
\midrule
squish & +0.092 & +0.073 & -- & -- \\
short-centre & +0.085 & +0.070 & 0.000 & 0.000 \\
letterbox & +0.092 & +0.069 & $-0.004$* & $-0.004$* \\
\bottomrule
\end{tabular}
\caption{Differences in AUC for the aspect-preserving presentations, with the
frozen crop-trained branches on \texttt{ntire\_val\_hard} and a paired
bootstrap of 2000 resamples. Every robust gain
has a 95\% interval that excludes zero. * marks a clean difference whose
interval excludes zero.}
\label{tab:aspect}
\end{table}

\section*{S3. Degradation after squish}

Under the degradation grid of the main text the degradation is applied
before squish, so squish's downscale could average it away. We registered a
second arm that squishes the clean image first and degrades the squished
view. The share of the rise with degradation level that survives is 0.86
[0.83, 0.89] on branch A and 0.55 [0.47, 0.62] on B. B's input is 224 pixels,
so squish downscales its images 2.3 times more than A's. The registered rule
gives A: coverage, B: unresolved. Degradation on the small squished view is
not the same strength as native degradation seen through the enlarged window,
so this arm removes the downscale's averaging without equalising what the two
arms see. The stronger-downscale test in the main text asks the same
question directly.

\section*{S4. DINOv2 test-time comparison on the proxies (exploratory, moved verbatim)}

The CLIP results led us to run the same test-time comparison on the DINOv2
proxies, which we had not registered. Shown the whole image, the DINOv2
crop-trained proxies gain +0.026 to +0.038 robust AUC, every interval clear
of zero, and the gain again grows with size (+0.075 above 752 pixels, +0.023
below 336). At proxy scale, training under squish adds to the test-time
change for both backbones: squish-trained proxies beat the crop-trained ones
shown the whole image by +0.048 to +0.057 robust and +0.022 to +0.032 clean
with DINOv2, and by +0.038 to +0.048 and +0.025 to +0.040 with CLIP, every
interval clear of zero. Test-time squish costs the crop-trained proxies 1 to
4 clean points, which training recovers.

\begin{table}[h]
\centering
\small
\begin{tabular}{llcc}
\toprule
backbone & model & robust & clean \\
\midrule
DINOv2 & crop-trained & 0.7994 & 0.9255 \\
       & crop-trained, squish at test & 0.8310 & 0.9087 \\
       & squish-trained & 0.8841 & 0.9357 \\
\midrule
CLIP & crop-trained & 0.7482 & 0.9655 \\
     & crop-trained, squish at test & 0.8485 & 0.9324 \\
     & squish-trained & 0.8924 & 0.9645 \\
\bottomrule
\end{tabular}
\caption{Proxy models (one tower, one epoch on a quarter of the training set)
on \texttt{ntire\_val\_hard}, mean of three seeds. The DINOv2 test-time row
was not registered.}
\end{table}

CLIP mechanism check (registered): the training gain is +0.206 above 752
pixels against +0.091 below 336. On \texttt{ntire\_val} the DINOv2 proxy robust
gain was +0.035 to +0.037.

\section*{S5. Earlier method changes (moved verbatim)}

Before finding the coverage effect we pre-registered three changes aimed
at the degradations, each with a written gate. Distilling degraded features
toward clean ones lost about one robust point on the proxy, and under a
control that gave both views the same window it still gained nothing
($-0.2$ points against a gate of $+1$).
Conditioning the classifier on a predicted degradation embedding through FiLM
tied the baseline (0.803 against 0.804). Weighting the fusion by predicted
degradation severity could not clear its gate of $+0.5$ points even with an
oracle (+0.4 points pooled). These are single-proxy results, but together with the probe they
suggest the representation was not the bottleneck.

\section*{S6. Masking}

We registered that masking the periphery of the squished image would cost
more than masking its centre. It failed on both branches: masking a central
quarter cost about twice as much as masking an outer rim of equal area (A
$-0.013$ [$-0.021$, $-0.004$], B $-0.012$ [$-0.020$, $-0.003$]). A 25\%-area
rim is a thin border, the emptiest part of a photo, so the test compared the
busiest quarter with the emptiest. Unregistered: squish with the centre
masked still beats the centre window (A +0.069, B +0.055).

\section*{S7. Locating the clean cost}

On Synthbuster, branch A changes by at most 0.007 on every generator. Branch
B loses in two unrelated places, DALL-E 2 ($-0.110$, 1024 pixels) and Firefly
($-0.040$, the whole bucket above 1100 pixels). The registered rule required
one location and found none. On logits z-scored against each model's own
reals, large fakes gain less than small ones (A $-5.1$, B $-10.7$ real-sd),
and TPR at 1\% FPR on Firefly falls (A $-0.062$, B $-0.259$). Synthbuster is
near ceiling for A, so a null AUC result says little. The test-split
comparison of 0.974 and 0.880 also changes final against averaged weights,
training round, and a 518 against 512 run setting.

Moved verbatim from the old main text: Split by native size, the real images
are not the cause: under crop their mean score drifts toward ``fake'' as size
grows, and under squish it is flat. On their clean images above 1100 pixels,
branch A trained under squish scores 1.000 against 0.989 for crop on
\texttt{ntire\_val} (871 images) and 0.978 against 0.981 on
\texttt{ntire\_val\_hard} (295 images).

\section*{S8. Superseded and unidentifiable tests}

\paragraph{Frequency matching.} We squished the image, resized it
back to native size, and took the centre window, so the window had squish's
frequency content at its own small field of view. It recovered none of the
gain (A $-0.009$, B $-0.014$ against the centre window), which the registered
rule read as coverage. Above 752 pixels the round trip leaves a window that
is mostly blur, unlike anything in training, and it also costs clean AUC, so
the design is confounded. Superseded by the decomposition in the main text.

\paragraph{Shrink ratio against field of view.} A regression of the
gain on both could not identify them, because they move together. No result.

\paragraph{Degradation families.} JPEG rows gain 2 to 5 times more
than others. An exploratory headroom check shows most of that excess is
headroom: it falls to 1.2 to 1.6 times once the room left to gain is counted.

\paragraph{Text corrections.} Ten wording fixes with numbers,
applied in the text.

\section*{S9. Degradation grid details}

Pooled level-averaged gain of squish over the native presentation, none /
light / medium / heavy: A $-0.0011$ / +0.0011 / +0.0075 / +0.0285, B +0.0043
/ +0.0076 / +0.0141 / +0.0271, B-Free $-0.1139$ / $-0.1080$ / $-0.0884$ /
$-0.0346$. The sign change on clean images is resolved for A only, in the
528 to 736 pixel bucket ($-0.0013$ [$-0.0027$, $-0.0003$]). A pre-stated
prediction that JPEG would rise most failed, and cannot be tested because the
levels are not matched in strength (pixel RMSE at heavy: JPEG 5.7, blur 11.1,
noise 23.6). On Synthbuster squish loses at every degradation level.

\end{document}
